\section{Missiooni eesmärgid}

\subsection{Esmane missioon}

Kõik CanSatid peavad täitma esmase missiooni, mis on ette nähtud satelliidi baassüsteemide funktsionaalsuse demonstreerimiseks: õhutemperatuuri ja -rõhu mõõtmine reaalajas ning edastamine iga sekund maapealsele sidejaamale.

\subsection{Teisene missioon}

Teise missiooni eesmärgi valivad meeskonnad ise. Meie meeskond valis õhuniiskuse ja CO$_2$-sisalduse mõõtmine SCD41 anduri kasutades.

\subsection{Mida on vaja eesmärkide saavutamiseks}

Eesmärkide saavutamiseks on vajalik kokkuhoidvat meeskonda, toetavat juhendajat, häid teadmisi ja motivatsiooni õppimise vastu ning loomulikult ressursse.

\section{Meeskonna korraldus ja rollid}

\begin{table}[H]
    \centering
    \noindent\setlength\tabcolsep{4pt}%
    \begin{tabularx}{\linewidth}{|l|l|l|>{\RaggedRight\arraybackslash}X|}
        \hline
        \textbf{nr} & \textbf{nimi} & \textbf{vanus} & \textbf{roll}\\
        \hline
        \hline
        1 & Artjom Jemeljanov & 18 & meeskonna liider, tööde jagaja, vastutab elektroonilise osa eest --– programmeerimine, signaali vastuvõtmine, andurite lisamine\\
        \hline
        2 & Anastassia Jurtšenko & 16 & vastutab tehnilise osa eest: langevari ja korpus, valmistab erinevate langevarjude mudelid erinevatest materjalidest ning erineva konstruktsioonidega\\
        \hline
        3 & Artjom Antonov & 18 & koostab aruanne, korjab vajalikud andmed, pildid, dokumendid\\
        \hline
        4 & Albert Shtele& 17 & vastutab CanSati andmete vastuvõtu eest, korpuse osade valmistamise eest ning vajaliku materjalide tellimise eest\\
        \hline
    \end{tabularx}
    \caption{Meeskonna liikmed}
    \label{tab:1}
\end{table}

Projekti alguses meeskonnas olid veel kaks liiget, kuid nad loobusid projektis osalemast ning meie jagasime nende ülesanded omavahel.